% docs/slides/preamble.tex -- 五份 deck 共享
\documentclass[aspectratio=169,11pt]{beamer}

% ==== 中文与字体 ====
\usepackage[UTF8,fontset=none]{ctex}
\usepackage{fontspec}

\IfFontExistsTF{PingFang SC}{%
  \setCJKmainfont{PingFang SC}[BoldFont=PingFang SC,ItalicFont=PingFang SC,AutoFakeSlant=0.15]
  \setCJKsansfont{PingFang SC}[ItalicFont=PingFang SC,AutoFakeSlant=0.15]
  \setCJKmonofont{PingFang SC}
}{%
  \IfFontExistsTF{Source Han Sans SC}{%
    \setCJKmainfont{Source Han Sans SC}
    \setCJKsansfont{Source Han Sans SC}
  }{%
    \IfFontExistsTF{Noto Sans CJK SC}{%
      \setCJKmainfont{Noto Sans CJK SC}
      \setCJKsansfont{Noto Sans CJK SC}
    }{%
      \setCJKmainfont{Heiti SC}
      \setCJKsansfont{Heiti SC}
    }
  }
}

\IfFontExistsTF{Helvetica Neue}{\setmainfont{Helvetica Neue}}{%
  \IfFontExistsTF{Inter}{\setmainfont{Inter}}{}}
\IfFontExistsTF{JetBrains Mono}{\setmonofont{JetBrains Mono}[Scale=0.85]}%
  {\setmonofont{Menlo}[Scale=0.85]}

% ==== 颜色 ====
\usepackage{xcolor}
\definecolor{cMain}{HTML}{1F4E79}
\definecolor{cAccent}{HTML}{E07A1F}
\definecolor{cGray}{HTML}{595959}
\definecolor{cMute}{HTML}{F4F4F4}
\definecolor{cOK}{HTML}{2E7D32}
\definecolor{cBad}{HTML}{B23A3A}
\definecolor{cWarn}{HTML}{C8A300}

% ==== 主题 ====
\usepackage{tcolorbox}
\usetheme{metropolis}
\metroset{progressbar=frametitle,numbering=fraction,block=fill}
\setbeamercolor{frametitle}{bg=cMain,fg=white}
\setbeamercolor{progress bar}{fg=cAccent,bg=cMute}
\setbeamercolor{alerted text}{fg=cAccent}
\setbeamercolor{title}{fg=cMain}

% metropolis 默认尝试 Fira Sans，未安装时 fallback 到 Latin Modern Sans。
% 这里在主题加载后显式覆盖，强制使用 macOS 上一定存在的 Helvetica Neue，
% 跟中文 PingFang 风格一致；找不到则按顺序尝试 Inter / Arial。
\IfFontExistsTF{Helvetica Neue}{%
  \setsansfont{Helvetica Neue}%
}{%
  \IfFontExistsTF{Inter}{\setsansfont{Inter}}{%
    \IfFontExistsTF{Arial}{\setsansfont{Arial}}{}}%
}

% ==== TikZ ====
\usepackage{tikz}
\usetikzlibrary{positioning,arrows.meta,fit,shapes.geometric,calc,%
  decorations.pathmorphing,backgrounds}

\tikzset{
  node-box/.style={draw=cMain,thick,rounded corners=2pt,
    minimum height=8mm,minimum width=18mm,align=center,font=\small},
  node-state/.style={draw=cMain,thick,circle,minimum size=14mm,
    align=center,font=\small},
  node-ok/.style={draw=cOK,thick,fill=cOK!10},
  node-warn/.style={draw=cWarn,thick,fill=cWarn!15},
  node-bad/.style={draw=cBad,thick,fill=cBad!10},
  % 色盲友好的状态机三态：蓝 (正常/Closed) / 橙 (告警/Open) / 灰 (中间/HalfOpen)
  % 避免 red-green 二元区分，改用 hue + lightness 双通道。
  node-state-normal/.style={draw=cMain,thick,fill=cMain!12},
  node-state-alert/.style={draw=cAccent,thick,fill=cAccent!18},
  node-state-transition/.style={draw=cGray,thick,fill=cGray!18},
  arrow/.style={-{Stealth[length=2mm]},thick,draw=cMain},
  arrow-async/.style={-{Stealth[length=2mm]},thick,dashed,draw=cGray},
  arrow-bad/.style={-{Stealth[length=2mm]},thick,draw=cBad},
  arrow-alert/.style={-{Stealth[length=2mm]},thick,draw=cAccent},
  lifeline/.style={dashed,thick,draw=cGray},
}

% ==== 代码 ====
\usepackage{listings}
\lstdefinelanguage{Go}{
  morekeywords={break,case,chan,const,continue,default,defer,else,
    fallthrough,for,func,go,goto,if,import,interface,map,package,range,
    return,select,struct,switch,type,var,nil,true,false,iota},
  morekeywords=[2]{string,int,int64,uint,uint64,bool,byte,error,context},
  sensitive=true,
  morecomment=[l]//,morecomment=[s]{/*}{*/},
  morestring=[b]",morestring=[b]`,
}
\lstdefinelanguage{LuaRedis}{
  morekeywords={and,break,do,else,elseif,end,false,for,function,if,in,
    local,nil,not,or,repeat,return,then,true,until,while},
  morekeywords=[2]{redis,call,KEYS,ARGV,tonumber,HGET,HSET,GET,SET,
    DECR,DECRBY,INCR,INCRBY,EXPIRE,PEXPIRE,ZADD,ZCARD,ZREMRANGEBYSCORE,
    EVAL,EVALSHA},
  sensitive=true,
  morecomment=[l]{--},
  morestring=[b]",morestring=[b]',
}

\lstdefinestyle{base}{
  basicstyle=\ttfamily\scriptsize,
  numbers=left,numberstyle=\tiny\color{cGray},
  keywordstyle=\color{cMain}\bfseries,
  keywordstyle=[2]\color{cAccent},
  stringstyle=\color{cOK},
  commentstyle=\color{cGray}\itshape,
  backgroundcolor=\color{cMute},
  frame=single,framerule=0pt,
  showstringspaces=false,
  breaklines=true,
  tabsize=2,
  columns=fullflexible,
}
\lstset{style=base}

% ==== 三线表与附录页码 ====
\usepackage{booktabs}
\usepackage{appendixnumberbeamer}

% ==== 页脚来源标注 ====
\newcommand{\srcnote}[1]{{\color{cGray}\scriptsize\itshape #1}}

% 关键论点框
\newcommand{\keypoint}[1]{%
  \begin{tcolorbox}[colback=cMute,colframe=cMain,boxrule=0.5pt,arc=2pt]%
  #1\end{tcolorbox}}


\title{订单生命周期：从下单到结清的 7 个业务节点}
\subtitle{C 端 / 商家 / 客服 / 运营 / SRE 各看到什么 — gomall v2.2 (business-first)}
\author{gomall · 业务侧 deck}
\date{}

\begin{document}

\maketitle

\begin{frame}{今日大纲}
\tableofcontents
\end{frame}

% ============================================================
\section{业务困局：一笔订单要回答 5 个角色的问题}
% ============================================================

\begin{frame}{订单状态不是 type 字段，是 5 个角色的协同语言}
\begin{tikzpicture}[node distance=5mm and 10mm,every node/.style={font=\scriptsize}]
  \node[node-box,fill=cMain!12] (ord) {\textbf{order.type}\\\tiny 一个 tinyint};
  \node[node-box,right=of ord,node-warn] (c) {C 端\\\tiny 我的货到哪了};
  \node[node-box,above=of c,node-warn] (m) {商家\\\tiny 我该发哪单};
  \node[node-box,below=of c,node-warn] (cs) {客服\\\tiny 怎么回复用户};
  \node[node-box,right=of m,node-warn] (ops) {运营\\\tiny GMV 漏斗哪截断};
  \node[node-box,right=of cs,node-warn] (sre) {SRE\\\tiny 关单链路有没有掉单};
  \draw[arrow] (ord) -- (c);
  \draw[arrow] (ord) -- (m);
  \draw[arrow] (ord) -- (cs);
  \draw[arrow] (ord) -- (ops);
  \draw[arrow] (ord) -- (sre);
\end{tikzpicture}
\vspace{2mm}
\begin{itemize}
  \item 一个 \texttt{type} 字段同时是\textbf{C 端跟踪页面文案}、\textbf{商家工作台筛选条件}、\textbf{客服 SOP 入口}、\textbf{运营漏斗节点}、\textbf{SRE 关单监控指标}。
  \item 状态切得不够细，5 个角色就要在工单 / IM / 群里口头补 — 这是\textbf{真正的业务债}。
  \item gomall 把 type 拆到 7 态，是给这 5 个角色一个\textbf{统一的事实表 (single source of truth)}。
\end{itemize}
\srcnote{consts/order.go:6--16 状态常量（7 态 + 拼团过渡态 8）；consts/order.go:19--28 OrderStateMap 提供中文名给前台 / 客服}
\end{frame}

\begin{frame}{3 态够吗：旧版本卡住的真实业务症状}
\begin{itemize}
  \item 旧版本只有 \texttt{UnPaid / Paid / Cancelled} 三态，C 端付了款只能看到\textbf{"已支付"}，但货还没发。
  \item \textbf{客服日均高频客诉 \#1}：\textit{我付了 3 天还没发货，你们是不是跑路了}。客服只能 IM 问商家，回话延迟在小时级。
  \item \textbf{商家工作台无筛选条件}：商家不知道哪些是\textbf{待发货}、哪些是\textbf{用户已收货}，每天靠人肉 sheet 维护。
  \item \textbf{运营漏斗失真}：付款 $\to$ 发货 $\to$ 收货三个节点合并成 "Paid"，无法算 \textbf{发货时效} / \textbf{签收时效} / \textbf{自动确认率}。
  \item \textbf{售后 / 退款无落点}：用户申请退款无独立态，业务方只能在 \texttt{Paid} 上加 flag，越加越脏。
\end{itemize}
\keypoint{业务侧的状态划分粒度 = \textbf{角色之间需要协同的最小事件}。3 态合 5 件事，必爆。}
\srcnote{旧 consts/order.go 仅 UnPaid/Paid/Cancelled 三态；docs/blog/15-order-close.md 是这次扩展的前置背景}
\end{frame}

\begin{frame}{gomall v2.2 的 7 态决策：业界对齐 + 业务诉求}
{\small
\begin{tabular}{@{}llll@{}}
\toprule
gomall & 中文 & 触发事件 & 主导角色 \\
\midrule
WaitPay (1)     & 待付款       & 用户下单成功                     & C 端 \\
WaitShip (2)    & 待发货       & 支付通道回调成功                 & 商家 \\
WaitReceive (4) & 已发货待收货 & 商家点"发货"                    & 物流 / 商家 \\
Completed (5)   & 已完成       & 用户确认 \textbf{或} 7d 自动     & 平台 cron \\
Closed (3)      & 已关闭       & 30min 未付 \textbf{或} 用户取消  & 关单链路 \\
Refunding (6)   & 退款中       & 用户发起退款                     & 客服 / 商家 \\
Refunded (7)    & 已退款       & 商家 / 运营批准                  & 资金 / 钱包 \\
\bottomrule
\end{tabular}}
\vspace{2mm}
\begin{itemize}
  \item 7 态 = 业界（淘宝 / 京东 / 拼多多 / Amazon）共有的\textbf{最低公约数}，做小要回头补、做大需要部分发货 / 多 SKU 子单。
  \item \textbf{数值兼容}：1/2/3 沿用旧值（存量数据零迁移），4-7 是新增；旧名 \texttt{UnPaid / Paid / Cancelled} 曾以 Deprecated 别名过渡，调用方迁移完成后\textbf{已删除}，现只剩 \texttt{OrderWaitPay} 等现行命名。
\end{itemize}
\srcnote{consts/order.go:6--16 数值定义；Deprecated 别名块已随调用方迁移完成删除}
\end{frame}

% ============================================================
\section{7 态 × 3 视角：C 端 / 商家 / 客服 各看到什么}
% ============================================================

\begin{frame}{帧 1 · WaitPay：C 端 / 商家 / 客服 各看到什么}
{\footnotesize
\begin{tabular}{@{}lp{3.5cm}p{3cm}p{3cm}@{}}
\toprule
角色 & 看到 & 能做 & 何时切走 \\
\midrule
C 端 & 订单卡片 \textit{"待付款"} + 倒计时 30:00 & 立即支付 / 取消 & 付完 $\to$ WaitShip \\
商家 & \textbf{不显示} & 无操作 & — \\
客服 & 倒计时进度条 + 业务码 70001 / 60002 可能 & 解释关单时长 / 重发支付 & 30min 后 $\to$ Closed \\
\bottomrule
\end{tabular}}
\vspace{2mm}
\begin{itemize}
  \item 业务关键：\textbf{30min 倒计时}是 C 端最敏感的体验点 — 太短伤"研究着买"用户，太长锁库存伤商家。
  \item 商家\textbf{故意看不到} WaitPay 订单 — 没付钱的不算"待发货工作量"，否则商家工作台数据失真。
  \item 客服 SOP：用户问"为什么 30min 关单"$\to$ 引用平台规则 + 引导重下；用户付款失败 $\to$ 排查支付通道。
\end{itemize}
\srcnote{repository/rabbitmq/order\_delay.go:22 OrderCancelDelay=30min；middleware/idempotency.go 60002 业务码}
\end{frame}

\begin{frame}{帧 2 · WaitShip：商家工作流的真正起点}
{\footnotesize
\begin{tabular}{@{}lp{3.5cm}p{3cm}p{3cm}@{}}
\toprule
角色 & 看到 & 能做 & 何时切走 \\
\midrule
C 端 & \textit{"待商家发货"} + 已付金额 & 申请退款 (未发) & 商家发货 $\to$ WaitReceive \\
商家 & 工作台\textbf{【待发货】tab} & 打单 / 发货 / 备注 & 点"发货"切到 WaitReceive \\
客服 & 商家发货时效 SLA & 催商家 / 同意退款 & — \\
\bottomrule
\end{tabular}}
\vspace{2mm}
\begin{itemize}
  \item \textbf{客服客诉 \#2}：\textit{我付款 3 天显示待发货}。客服界面能看到\textbf{付款时间戳 + 商家最近登录}，秒回 SOP。
  \item 业务承诺：商家 \textbf{48h 未发货} $\to$ 用户可一键申请退款（路线图：超时自动拉到 Refunding）。
  \item \textbf{平台 P0 指标}：\texttt{付款 $\to$ 发货} 的 p95 时延 = 平台对商家的硬约束。
\end{itemize}
\srcnote{internal/order/shipping.go:36--67 ShipOrder；events/events.go OrderShipped 事件携带 tracking}
\end{frame}

\begin{frame}{帧 3 · WaitReceive：平台 7d 自动确认的资金钟}
{\footnotesize
\begin{tabular}{@{}lp{3.5cm}p{3cm}p{3cm}@{}}
\toprule
角色 & 看到 & 能做 & 何时切走 \\
\midrule
C 端 & \textit{"已发货"} + 物流单号 & 确认收货 / 申请退款 & 点确认 / 7d 到 $\to$ Completed \\
商家 & 工作台\textbf{【已发货】tab} & 跟物流 / 处理拒收 & — \\
客服 & 物流停滞天数 + 区域 & 协调物流 / 同意全额退 & — \\
\bottomrule
\end{tabular}}
\vspace{2mm}
\begin{itemize}
  \item 业务事实：\textbf{绝大多数用户不主动点确认收货} — 收到货之后没人会专门回 App 点一下。
  \item 平台 cron 每 6 小时扫一次 \texttt{updated\_at > 7d} 的 WaitReceive 订单 $\to$ 自动推进到 Completed。
  \item 这个动作是\textbf{资金结算给商家的触发器}：商家货款 = Completed 后 T+1 入账（路线图）。
\end{itemize}
\srcnote{internal/order/task.go:61--94 RunAutoConfirmReceive；initialize/cron.go:47 注册 0 0 */6}
\end{frame}

\begin{frame}{帧 4 · Completed / Closed / Refunding / Refunded：三终态 + 一中间}
{\footnotesize
\begin{tabular}{@{}lp{2.8cm}p{2.8cm}p{2.5cm}@{}}
\toprule
状态 & C 端体验 & 商家 / 平台 & 资金 \\
\midrule
Completed & 评价开窗 / 售后期内 & 货款待结算 & 进商家结算池 \\
Closed    & "已关闭" + 重下入口 & \textbf{不出现} & 库存预占已释放 \\
Refunding & "退款审核中" + 进度 & 客服 / 商家审核 & 货款\textbf{冻结} \\
Refunded  & "退款已到账 / 进行中" & 已扣减 GMV & 退原路 / 站内余额 \\
\bottomrule
\end{tabular}}
\vspace{2mm}
\begin{itemize}
  \item \textbf{Refunding 是唯一可双向出边的中间态}：批准 $\to$ Refunded，驳回 $\to$ Completed（评价重开 / 售后工单关）。
  \item \textbf{Closed} 不进商家工作台：未付款的不算业务量；C 端可看到，方便引导"重新下单"挽回 GMV。
  \item \textbf{Completed} 不是死刑：售后窗口期内仍可发起退款 (\texttt{Completed $\to$ Refunding} 这条边存在)。
\end{itemize}
\srcnote{internal/order/state.go:21--28 转换表；internal/refund/service.go:47--170 三个退款方法}
\end{frame}

\begin{frame}{完整状态机：用户 / 商家 / 平台 各自的边}
\begin{center}
\begin{tikzpicture}[node distance=14mm and 16mm,every node/.style={font=\scriptsize}]
  \node[node-box,fill=cMain!12] (wp) {WaitPay\\\tiny (1)};
  \node[node-box,fill=cMain!12,right=of wp] (ws) {WaitShip\\\tiny (2)};
  \node[node-box,fill=cMain!12,right=of ws] (wr) {WaitReceive\\\tiny (4)};
  \node[node-box,fill=cMain!12,right=of wr] (cp) {Completed\\\tiny (5)};

  \node[node-box,fill=cGray!18,below=of wp] (cl) {Closed\\\tiny (3) 终态};
  \node[node-box,fill=cAccent!18,below=of wr] (rg) {Refunding\\\tiny (6)};
  \node[node-box,fill=cGray!18,below=of cp] (rd) {Refunded\\\tiny (7) 终态};

  \draw[arrow] (wp) -- node[above,font=\tiny]{用户付款}(ws);
  \draw[arrow] (ws) -- node[above,font=\tiny]{商家发货}(wr);
  \draw[arrow] (wr) -- node[above,font=\tiny]{用户确认/7d}(cp);

  \draw[arrow-bad] (wp) -- node[left,font=\tiny]{30min/取消}(cl);
  \draw[arrow-alert] (ws.south east) -- node[above,sloped,font=\tiny\color{cAccent}]{用户退款}(rg.north west);
  \draw[arrow-alert] (wr) -- node[right,font=\tiny\color{cAccent}]{用户退款}(rg);
  \draw[arrow-alert] (cp.south west) -- node[above,sloped,font=\tiny\color{cAccent}]{用户退款}(rg.north east);

  \draw[arrow] (rg) -- node[above,font=\tiny]{商家批}(rd);
  \draw[arrow,dashed] (rg.north east) to[bend right=15] node[right,pos=0.4,font=\tiny]{商家驳回}(cp.south);
\end{tikzpicture}
\end{center}
\vspace{2mm}
\begin{itemize}
  \item \textbf{蓝边} = 正向流转（业务期望）；\textbf{橙边} = 用户触发退款；\textbf{灰边} = 关单。
  \item 每条边都有\textbf{触发主语}（用户 / 商家 / cron）：业务上不允许商家替用户确认收货（防薅平台羊毛）。
\end{itemize}
\srcnote{internal/order/state.go:21--28 转换表；internal/order/shipping.go:36--111 + refund/service.go:47--170}
\end{frame}

% ============================================================
\section{关单博弈：30min 倒计时背后的业务取舍}
% ============================================================

\begin{frame}{30min 倒计时：体验 vs 锁库存的钢丝绳}
\begin{tikzpicture}[node distance=6mm and 14mm]
  \node[node-box,node-warn] (short) {太短 \textit{(<10min)}\\\scriptsize 用户研究着买被关掉};
  \node[node-box,node-ok,right=of short] (mid) {\textbf{30min}\\\scriptsize 业界共识};
  \node[node-box,node-bad,right=of mid] (long) {太长 \textit{(>2h)}\\\scriptsize 库存锁死，商家失血};
  \node[node-box,below=of short] (c) {C 端体验};
  \node[node-box,below=of mid] (b) {业务取舍点};
  \node[node-box,below=of long] (m) {商家利益};
  \draw[arrow] (short) -- (mid);
  \draw[arrow] (mid) -- (long);
  \draw[arrow,dashed] (c) -- (short);
  \draw[arrow,dashed] (m) -- (long);
\end{tikzpicture}
\vspace{2mm}
\begin{itemize}
  \item \textbf{业务事实}：用户加购到付款的中位时长 5-8min（业界公开数据），\textbf{30min} = 给 95\%+ 用户足够时间 + 给商家可控的库存锁定窗口。
  \item \textbf{平台四方对齐}：淘宝 30min / 京东 30min / 拼多多 30min / Amazon 30min — gomall 对齐业界。
  \item \textbf{业务规则}：超时关单同时\textbf{释放 Redis reserved 预占库存}，让下一批用户能正常下单。
\end{itemize}
\srcnote{repository/rabbitmq/order\_delay.go:22 OrderCancelDelay=30min；internal/order/cancel.go:61 释放 reserved}
\end{frame}

\begin{frame}{为什么单一保险不够：RMQ 漏 + Cron 慢}
\begin{center}
\begin{tikzpicture}[node distance=12mm and 14mm,every node/.style={font=\scriptsize}]
  \node[node-box,fill=cMain!12,minimum width=22mm,minimum height=12mm] (order) {下单};
  \node[node-box,right=of order,fill=cMain!12,minimum width=22mm,minimum height=12mm] (rmq) {RMQ TTL\\\tiny 精确到秒};
  \node[node-box,right=of rmq,node-ok,minimum width=22mm,minimum height=12mm] (close1) {关单};
  \node[node-box,below=of order,fill=cMain!12,minimum width=22mm,minimum height=12mm] (db) {DB type=1};
  \node[node-box,right=of db,fill=cMain!12,minimum width=26mm,minimum height=12mm] (cron) {Cron 5min 扫\\\tiny WaitPay + timeout};
  \node[node-box,right=of cron,node-ok,minimum width=22mm,minimum height=12mm] (close2) {关单};
  \draw[arrow] (order) -- (rmq);
  \draw[arrow] (rmq) -- (close1);
  \draw[arrow] (order) -- (db);
  \draw[arrow,dashed] (db) -- (cron);
  \draw[arrow,dashed] (cron) -- (close2);
\end{tikzpicture}
\end{center}
\vspace{2mm}
\begin{itemize}
  \item \textbf{单 RMQ 不够}：RMQ 重启 / 消息丢失 / 消费者挂掉 = 漏关一单 = \textbf{库存永久锁死}。
  \item \textbf{单 Cron 不够}：5min 跑一次 $\to$ 实时性差，用户付款失败后看着"待付款"挂十几分钟。
  \item \textbf{双保险叙事}：RMQ 走\textbf{快线}（精确到秒），Cron 走\textbf{慢线}（兜底）— 漏一个不漏单。
  \item 业务承诺：\textbf{不允许漏关一单}。两条路共享同一个 \texttt{CancelUnpaidOrder} 入口。
\end{itemize}
\srcnote{repository/rabbitmq/order\_delay.go:59--72 PublishOrderCancelDelay；initialize/cron.go:21 cron 注册；internal/order/cancel.go:19--67 共享入口}
\end{frame}

\begin{frame}{时序图：下单落库 + 超时延迟关单的完整交互}
\centering
\resizebox{!}{0.50\textheight}{%
\begin{sequencediagram}
  \newthread{u}{用户}
  \newinst[1]{h}{Order Handler}
  \newinst[1]{s}{库存 Redis Lua}
  \newinst[1]{d}{订单 DAO}
  \newinst[1]{q}{RabbitMQ 延迟队列/DLX}
  \newinst[1]{c}{延迟关单消费者}

  \begin{call}{u}{下单 create}{h}{待支付}
    \begin{call}{h}{ReserveStock}{s}{available$\to$reserved}
    \end{call}
    \begin{call}{h}{snowflake 单号 + 落库 + outbox（同事务）}{d}{type=待支付}
    \end{call}
  \end{call}

  \mess[1]{h}{ZAdd + PublishOrderDelay（TTL 30min，异步）}{q}

  \mess[1]{q}{TTL 到期投递（异步）}{c}
  \begin{call}{c}{若仍未支付：关闭 + 释放 reserved}{d}{type=已关闭}
  \end{call}
  \mess[1]{c}{失败按投递次数进 DLX/DLQ（不无限 requeue）}{q}
\end{sequencediagram}
}

\vspace{0.5mm}
{\scriptsize
\begin{itemize}\itemsep0pt
  \item \textbf{下单同步段}：预扣库存 + snowflake 单号 + 落库 + 写 outbox 一个事务收口，对外只产出"待支付"。
  \item \textbf{延迟关单段}：RabbitMQ 走 \textbf{TTL 延迟}快线，到期投关单消费者；仍未支付则置"已关闭"并释放 reserved 预占。
  \item \textbf{双保险 + 毒丸隔离}：\textbf{RMQ TTL 延迟 + Cron 兜底}共享同一关单入口；坏消息按投递次数进 \textbf{DLQ}，不无限 requeue。
\end{itemize}}
\srcnote{internal/order/service.go:119/127/162--192；order\_delay.go:22/59--72 TTL；cancel.go:19--67 关单+释放；poison.go RouteToDLQ}
\end{frame}

\begin{frame}[fragile]{毒丸消息：失败关单无限 requeue 把延迟队列打爆}
\begin{lstlisting}[language=Go,firstnumber=87,
  caption={repository/rabbitmq/order\_delay.go:87--98 失败消息分流},captionpos=b]
if err := handler(orderNum); err != nil {
    util.LogrusObj.Errorf("处理延迟关单失败 orderNum=%d err=%v\n", orderNum, err)
    // order.dead.queue 未配 DLX，requeue 会立即回到队首重投。
    // 仅允许首次失败重投一次；再次失败（Redelivered）转独立死信队列，
    // 避免持续失败的毒丸消息无界回灌、占满消费者。
    if d.Redelivered {
        RouteToDLQ(d, orderDeadQueue, orderDeadRouting, false)
        return
    }
    _ = d.Nack(false, true)
    return
}
\end{lstlisting}
\begin{itemize}\itemsep1pt
  \item \textbf{修前}：handler 失败一律 \texttt{Nack(requeue=true)} —— 一条 \texttt{CloseOrderWithCheck} 反复失败的订单（如关联 outbox 写库持续报错）回到队首立刻重投，形成\textbf{热循环}。
  \item \textbf{业务风险}：单条坏消息 100\%\,CPU 空转重投，把消费者吃满；后面排队的\textbf{正常超时订单被饿死}，30min 关单 SLA 失守、库存预占释放被卡住。
  \item \textbf{修后}：首投失败仍重试一次（吸收瞬时抖动）；判定 \texttt{Redelivered} 即超阈值，\texttt{RouteToDLQ} 显式投独立死信队列并 Ack，原队列\textbf{不再回灌}。
  \item 死信里的坏单留给运维人工补偿，正常关单流量不被一条毒丸拖垮。
\end{itemize}
\srcnote{repository/rabbitmq/order\_delay.go:87--98 失败分流；poison.go:85--87 ExceededDeliveryLimit 阈值；poison.go:93--115 RouteToDLQ 投死信并 Ack}
\end{frame}

\begin{frame}[fragile]{已修复复盘：cron 表达式陷阱与现行写法}
\begin{lstlisting}[language=Go,firstnumber=19,
  caption={initialize/cron.go:19--28 现行写法 @every 5m（左侧为陷阱对照）},captionpos=b]
c := cron.New(cron.WithSeconds())
orderService := new(order.OrderTaskService)
// 现行：@every 5m，每 5min 整体触发一次（语义明确、不踩秒位陷阱）
// 对照陷阱：" * */5 * * * * " 在 WithSeconds 下 = 秒位每秒 + 分位每 5min
//          = 实际 60x/5min，是这类六段表达式最常见的误写
_, err := c.AddFunc("@every 5m", func() {
    defer func() {
        if r := recover(); r != nil {
            util.LogrusObj.Errorf("Cron 任务发生 Panic: %v", r)
        }
    }()
    orderService.RunOrderTimeoutCheck()
})
\end{lstlisting}
\begin{itemize}
  \item \textbf{陷阱本身}：\texttt{* */5 * * * *} 在 \texttt{WithSeconds} 下读作\textbf{秒位 = 任意}（每秒）+ 分位每 5 min — 行业里最常踩的 \textbf{60×/5min} 误写。
  \item \textbf{现行做法}：用 \texttt{@every 5m} 取代六段表达式，语义一眼可读，从源头避开秒位歧义。
  \item \textbf{若误写会有的代价}：DB 每分钟扫一次 \texttt{WHERE type=1}，6M 行 order 表单峰压力 60×；outbox \texttt{order.cancelled} 写入量虚高让 SRE 误判关单暴涨。
  \item \textbf{为什么即便误写也不会超卖}：\texttt{CloseOrderWithCheck} 的 \texttt{WHERE type=1} 条件 UPDATE 天然幂等，重复扫只是 \texttt{RowsAffected=0} 空跑 — 正确性靠幂等闸兜住，频率正确性靠表达式选型。
\end{itemize}
\srcnote{initialize/cron.go:21 现行 @every 5m；internal/order/cancel.go:32--38 CloseOrderWithCheck 幂等兜底}
\end{frame}

\begin{frame}[allowframebreaks]{cron 陷阱复盘：误写代价与现行加固}
\begin{center}
\begin{tikzpicture}[node distance=5mm and 12mm,every node/.style={font=\scriptsize}]
  \node[node-box,node-bad] (db) {误写代价：DB 压力 60×\\\tiny SELECT 全表扫};
  \node[node-box,node-bad,right=of db] (log) {误写代价：日志洪水\\\tiny LogrusObj 60×};
  \node[node-box,node-bad,below=of db] (mon) {误写代价：监控失真\\\tiny outbox 写入计数};
  \node[node-box,node-warn,right=of mon] (risk) {误写代价：潜在误关\\\tiny 边界并发};
  \node[node-box,node-ok,right=of log] (fix) {现行：\\\tiny "@every 5m"};
\end{tikzpicture}
\end{center}
\footnotesize
\begin{itemize}\itemsep0pt
  \item \textbf{当初的可见症状}：dev 环境 DB CPU 周期性飙红 + logrus 日志量异常 $\to$ 引出排查并修复。
  \item \textbf{为什么始终没造成事故}：幂等兜底 + 关单不可逆 + 释放 reserved 也是 \texttt{IF reserved>=N THEN ...} 的 Lua 幂等。
  \item \textbf{现行加固}：表达式改用 \texttt{@every 5m}（每 5min 整体触发 1 次），从写法上消除秒位歧义。
  \item \textbf{为什么 deck 要保留这段复盘}：把"踩过的坑 + 怎么收口"讲清楚 — \textbf{在面试 / CR 中比"假装全部一帆风顺"更值钱}。
\end{itemize}
\keypoint{\footnotesize 幂等设计的隐藏 ROI：它不仅防止业务超卖，也\textbf{吸收了重启重跑、网络重试乃至当初的表达式误写} — 但频率正确性仍须靠表达式选型本身保证。}
\srcnote{initialize/cron.go:21 现行 @every 5m；internal/order/cancel.go:32 + internal/order/repo.go:198 兜底链}
\end{frame}

\begin{frame}[fragile]{CancelUnpaidOrder：4 个调用方共享一个幂等入口}
\begin{lstlisting}[language=Go,firstnumber=19,
  caption={internal/order/cancel.go:19--67 共享入口（节选）},captionpos=b]
func CancelUnpaidOrder(orderNum uint64) error {
    ctx := context.Background()
    baseDao := NewOrderDao(ctx)
    order, err := baseDao.GetOrderByOrderNum(orderNum)
    if err != nil { return err }
    if order == nil || order.ID == 0 { return nil }

    var closed bool
    err = baseDao.DB.Transaction(func(tx *gorm.DB) error {
        ok, err := NewOrderDaoByDB(tx).CloseOrderWithCheck(orderNum)
        if err != nil { return err }
        if !ok { return nil }  // 已经被别的路径关掉了，no-op
        closed = true
        return outbox.NewOutboxDaoByDB(tx).Insert(
            "order", "OrderCancelled", "order.cancelled", order.ID,
            events.OrderCancelled{ /* ... */ Reason: "timeout",
                PromoRuleID:        order.PromoRuleID,
                PromoDiscountCents: order.PromoDiscountCents})
    })
    if err != nil { return err }
    if !closed { return nil }
    if relErr := cache.ReleaseReservation(ctx,
        order.ProductID, int64(order.Num)); relErr != nil {
        util.LogrusObj.Errorf("release on cancel failed: %v", relErr)
    }
    return nil // 满减预算退还由 promo 侧消费事件异步完成
}
\end{lstlisting}
\begin{itemize}
  \item \textbf{4 个调用方}：RMQ 消费者 / Cron 兜底 / 客服后台 / 用户主动取消 — 共享同一个入口、同一套兜底。
  \item \texttt{closed=false} 直接 no-op，避免重复释放 reserved、避免重复写 outbox。
  \item \textbf{满减预算不在这里同步退}：\texttt{order.cancelled} 事件载荷自带 \texttt{promo\_rule\_id / promo\_discount\_cents}，promo 侧 \texttt{promo.release} 队列消费后退预算，\texttt{promo\_release} 台账（OrderID 唯一索引）吸收重复投递。
\end{itemize}
\srcnote{internal/order/cancel.go:19--67；CloseOrderWithCheck WHERE type=1 兜底幂等；internal/promo/consumer.go:25--48 消费侧}
\end{frame}

% ============================================================
\section{客服话术帧：业务码 + SOP 一一对照}
% ============================================================

\begin{frame}{客诉 \#1：30min 自动关单 — 客服怎么回}
\begin{itemize}
  \item \textbf{用户原话}：\textit{我下单 30 分钟没付款，订单怎么自动取消了？我不知道还有时间限制。}
  \item \textbf{客服 SOP}：
  \begin{itemize}
    \item 1. 道歉 + 解释平台规则（统一文案库）：\textit{为了保证库存周转，所有订单付款窗口为 30 分钟。}
    \item 2. 引导重新下单 + 给 \textbf{5 元代金券}（路线图：客服权限发券）补偿。
    \item 3. 后台核对：\texttt{order.type=3 + reason=timeout + cancel\_time} 在 outbox \texttt{order.cancelled} 事件里能查到。
  \end{itemize}
  \item \textbf{业务码不暴露给客服}：客服界面显示\textbf{中文 "已超时关闭"}，不显示 \texttt{type=3}。
  \item 留痕：客诉对话 + 补偿券发放都写工单，运营复盘超时关单率。
\end{itemize}
\srcnote{repository/rabbitmq/order\_delay.go:22 30min 平台规则；internal/order/cancel.go:48 reason=timeout 事件标记}
\end{frame}

\begin{frame}{客诉 \#2：付了 3 天没发货 — 客服怎么回}
\begin{itemize}
  \item \textbf{用户原话}：\textit{我付了 3 天了订单还在"待发货"，是不是你们跑路了？我要退款！}
  \item \textbf{客服 SOP}：
  \begin{itemize}
    \item 1. 查商家工作台\textbf{【待发货】tab 待处理数} + 商家最近登录时间。
    \item 2. 商家 24h 内有活动 $\to$ IM 催商家，挂\textbf{客服跟进工单}（4h 复查）。
    \item 3. 商家 48h 无响应 $\to$ 客服\textbf{代发起退款}（走 \texttt{RequestRefund} 调用，from=WaitShip 合法）。
    \item 4. 客服 6h 内\textbf{先于商家批准}退款（避免用户二次客诉）。
  \end{itemize}
  \item \textbf{平台承诺}：商家 48h 未发货，用户一键退款（路线图）；现版本靠客服主动 trigger。
\end{itemize}
\srcnote{internal/refund/service.go:37--41 refundAllowedFrom 含 WaitShip；internal/refund/service.go:47--87 RequestRefund 主流程}
\end{frame}

\begin{frame}{客诉 \#3：点了确认收货能撤回吗 — 客服怎么回}
\begin{itemize}
  \item \textbf{用户原话}：\textit{我刚点确认收货，发现东西有问题，能撤回吗？}
  \item \textbf{客服 SOP}：
  \begin{itemize}
    \item 1. \textbf{不能撤回} — 状态机不允许 Completed $\to$ WaitReceive 回退（业务规则：避免薅羊毛）。
    \item 2. 但是 Completed 不是死刑 — \textbf{售后期内仍可发起退款}（\texttt{Completed $\to$ Refunding} 这条边存在）。
    \item 3. 引导用户走\textbf{售后申请}：\texttt{POST /orders/refund/request}，理由选"商品质量问题"。
    \item 4. 客服后台一键\textbf{加急工单}（路线图：售后工单服务落地后挂"加急" flag）。
  \end{itemize}
  \item \textbf{为什么不允许撤回}：用户撤回 = 资金回退商家结算 = T+1 入账的钱已经在路上 $\to$ \textbf{资金链反推不可行}。
\end{itemize}
\srcnote{internal/order/state.go:25 Completed 出边只有 Refunding；internal/refund/service.go:37--41 from 含 Completed}
\end{frame}

\begin{frame}{客诉 \#4：用户说没收到货 — 客服怎么回}
\begin{itemize}
  \item \textbf{用户原话}：\textit{我看物流显示签收了，但我没收到货！}
  \item \textbf{客服 SOP}：
  \begin{itemize}
    \item 1. 查\textbf{物流单详情}（路线图：logistics\_order 子表）+ 签收人 / 签收时间 / 签收地址。
    \item 2. 引导用户\textbf{联系小区 / 代收点 / 家人} 24h 内回查。
    \item 3. 仍找不到货 $\to$ 客服代发起\textbf{全额退款 + 物流商索赔}流程（路线图：物流 webhook 联动）。
    \item 4. 订单仍在 \texttt{WaitReceive}（用户没点确认）$\to$ 直接 RequestRefund。
    \item 5. 订单已到 \texttt{Completed}（7d 自动）$\to$ 走售后期申诉，靠 \texttt{Completed $\to$ Refunding} 边推进。
  \end{itemize}
\end{itemize}
\srcnote{internal/refund/service.go:37--41 from 列表覆盖 WaitReceive / Completed；事件 events/events.go OrderShipped 带 tracking}
\end{frame}

\begin{frame}{客诉 \#5：抢购 / 支付出现 70001 / 60002 — 客服怎么回}
\begin{itemize}
  \item \textbf{70001 限流}（\texttt{ErrRateLimitExceeded}）：\textit{您操作过于频繁，请稍后再试。}
  \begin{itemize}
    \item 客服话术：秒杀场景 1s 内 3 次以上点击会触发；引导用户\textbf{等 1 秒再点}。
    \item 不要解释"滑动窗口算法"，用户只关心"我是不是被封号了" — 答案是\textbf{没有，仅限频}。
  \end{itemize}
  \item \textbf{60002 幂等处理中}（\texttt{ErrIdempotencyInProgress}）：\textit{订单正在处理，请勿重复提交。}
  \begin{itemize}
    \item 用户在网络抖动时连点了 N 次 \texttt{/orders/create}，第二次起返回 60002。
    \item 客服话术：\textit{您的订单已提交，请回首页"我的订单"刷新查看。}
  \end{itemize}
  \item \textbf{70002 熔断打开}（\texttt{ErrCircuitOpen}）：\textit{当前业务繁忙，请稍后再试。}
  \begin{itemize}
    \item 下游服务大面积失败被熔断，客服关注\textbf{波及范围} $\to$ 转 SRE。
  \end{itemize}
\end{itemize}
\srcnote{pkg/e/code.go:57--58 70001 / 70002；pkg/e/code.go:54 60002 幂等}
\end{frame}

% ============================================================
\section{业务承诺总表与 SLO}
% ============================================================

\begin{frame}{业务承诺总表：7 条边 7 个业务事件}
{\small
\begin{tabular}{@{}p{3.5cm}p{3cm}p{3.5cm}@{}}
\toprule
业务承诺 & 状态边 & 触发主语 \\
\midrule
30min 未付自动关单     & WaitPay $\to$ Closed       & RMQ + Cron \\
商家发货切到履约       & WaitShip $\to$ WaitReceive & 商家 \\
7d 自动确认收货        & WaitReceive $\to$ Completed & 平台 cron \\
售后期内可申请退款     & WaitShip/Receive/Completed $\to$ Refunding & 用户 / 客服 \\
售后通过才真退款       & Refunding $\to$ Refunded   & 商家 / 运营 \\
售后驳回回到完成       & Refunding $\to$ Completed  & 商家 / 运营 \\
终态不可逆            & Closed / Refunded      & — \\
\bottomrule
\end{tabular}}
\vspace{2mm}
\begin{itemize}
  \item 7 条承诺 = 7 态状态机的核心；每条对应一条 SQL（条件 UPDATE）+ 一个 outbox 事件 + 一份单测。
  \item \textbf{业务确定性的来源}：每条边的"主语"是固定的 — 用户不能替商家发货，商家不能替用户确认收货。
\end{itemize}
\srcnote{internal/order/state.go:21--28；internal/order/shipping.go + refund/service.go + order/task.go 四件套}
\end{frame}

\begin{frame}{订单链路 SLO：从 README 业务承诺到压测对照}
{\small
\begin{tabular}{@{}p{4cm}p{2.5cm}p{2.5cm}@{}}
\toprule
业务节点 & 业务承诺 & 实测对照 \\
\midrule
下单（P0）             & 99.95\% / p99<500ms & \texttt{/orders/create} \textbf{50,319 RPS} / p95 2.33ms \\
订单列表（P1）         & 99.9\% / p99<200ms  & \texttt{/orders/list} \textbf{58,406 RPS} / p95 5ms \\
关单延迟（业务 SLA）   & ≤ 30min + 5min      & RMQ + Cron 双保险 \\
确认收货延迟           & ≤ 7d + 6h           & cron 6h 一跑 \\
退款发起 $\to$ 资金到账    & 1-15d（原路）       & saga 异步（钱包服务消费）\\
\bottomrule
\end{tabular}}
\vspace{2mm}
\begin{itemize}
  \item \textbf{宁可下游慢、不能上游 429}：下单 / 列表都\textbf{不挂限流}，靠缓存 + 异步削峰扛峰值。
  \item \textbf{留 5-10× 余量给 GC / 网络抖动}：实测 RPS × 10 = 业务平时峰值；压测验证容量 ceiling。
\end{itemize}
\srcnote{stressTest/REPORT.md 第 4 节 755K 请求 1 单；REPORT.md 第 2 节订单列表 7000× 提升}
\end{frame}

\begin{frame}{业务可观测指标：5 个状态切换的 metrics}
\begin{itemize}
  \item \textbf{每条边都打 metric}：\texttt{order\_state\_transition\_total\{from,to,trigger\}} —— 由 outbox publisher 消费事件落 Prometheus。
  \item \textbf{运营看板核心面板}：
  \begin{itemize}
    \item GMV 漏斗：WaitPay 创建 $\to$ WaitShip 转化率 $\to$ WaitReceive 转化率 $\to$ Completed 转化率
    \item 关单率：\texttt{Closed / (Closed + WaitShip)} 按小时切分（异常 $\to$ 支付通道 / 反欺诈警报）
    \item 发货 p95 时延：\texttt{WaitShip.updated\_at - WaitShip.created\_at}
    \item 自动确认占比：\texttt{Completed.Auto=true / Completed.total}（应 >80\% 业内基线）
    \item 退款率：\texttt{Refunded / Completed}（>3\% 触发告警，路线图）
  \end{itemize}
  \item \textbf{SRE 告警}：关单 RPS 突增（误关风险）/ 自动确认率突降（cron 失活）/ 退款率突增（商品问题 / 黑产）。
\end{itemize}
\srcnote{service/events/events.go 各事件 payload；outbox publisher 路线图接 Prometheus}
\end{frame}

% ============================================================
\section{状态机实现：3 层兜底（业务规则 / SQL / 事件）}
% ============================================================

\begin{frame}[fragile]{第 1 层：状态机表是业务规则，不是数据结构}
\begin{lstlisting}[language=Go,firstnumber=21,
  caption={internal/order/state.go:21--28 转换表},captionpos=b]
var orderStateTransitions = map[uint][]uint{
    consts.OrderWaitPay:     {consts.OrderWaitShip, consts.OrderClosed},
    consts.OrderWaitShip:    {consts.OrderWaitReceive, consts.OrderRefunding},
    consts.OrderWaitReceive: {consts.OrderCompleted, consts.OrderRefunding},
    consts.OrderCompleted:   {consts.OrderRefunding},
    consts.OrderRefunding:   {consts.OrderRefunded, consts.OrderCompleted},
    consts.OrderWaitGroup:   {consts.OrderWaitShip, consts.OrderClosed}, // 拼团过渡态 (deck 14)
}
\end{lstlisting}
\begin{itemize}
  \item \textbf{key} = 当前状态，\textbf{value} = 允许去的下一态集合 — 不在表里就是\textbf{业务上非法}。
  \item \textbf{Closed} 和 \texttt{Refunded} 不在 key 集合 $\to$ \texttt{IsTerminalOrderState} 直接 \texttt{true}。
  \item \textbf{为什么放 service 不放 model}：状态机是业务规则，DB 只兜底"原子推进"，业务正确性靠这张表 + 单测。
  \item 17 例单测覆盖：合法 / 非法 / 终态 / 未知输入。
\end{itemize}
\srcnote{internal/order/state.go:21--28 表；:36--47 CanTransition；:58--61 IsTerminal}
\end{frame}

\begin{frame}[fragile]{第 2 层：DAO 条件 UPDATE 是并发的最后一道闸}
\begin{lstlisting}[language=Go,firstnumber=215,
  caption={internal/order/repo.go:215--223 ShipOrder DAO},captionpos=b]
func (d *OrderDao) ShipOrder(orderNum uint64) (bool, error) {
    res := d.DB.Model(&Order{}).
        Where("order_num=? AND type=?", orderNum, consts.OrderWaitShip).
        Update("type", consts.OrderWaitReceive)
    if res.Error != nil { return false, res.Error }
    return res.RowsAffected > 0, nil
}
\end{lstlisting}
\begin{itemize}
  \item \texttt{WHERE type=WaitShip} 兜住所有并发：第二次重发只会 \texttt{RowsAffected=0}，service 层返回 \texttt{ErrInvalidOrderStateTransition}。
  \item 这条 SQL 是状态机推进的\textbf{原子保证}：service 层就算被绕过，DB 仍能拒绝非法推进。
  \item 退款的 \texttt{WHERE type IN (?)} 同理 — 多 from 列表也能在 SQL 层一次拦截。
\end{itemize}
\srcnote{internal/order/repo.go:215--223 ShipOrder；:227--235 ConfirmReceive；:239--250 RequestRefund 多 from}
\end{frame}

\begin{frame}[allowframebreaks,fragile]{第 3 层：outbox 事件解耦下游服务}
\begin{lstlisting}[language=Go,firstnumber=36,
  caption={internal/order/shipping.go:36--67 ShipOrder 主体（节选）},captionpos=b]
func (s *ShippingSrv) ShipOrder(ctx context.Context, orderNum uint64,
    trackingNo, carrier string) error {
    // ... 校验状态 + 取订单 ...
    return baseDao.DB.Transaction(func(tx *gorm.DB) error {
        ok, err := NewOrderDaoByDB(tx).ShipOrder(orderNum)
        if err != nil { return err }
        if !ok { return ErrInvalidOrderStateTransition }
        return outbox.NewOutboxDaoByDB(tx).Insert(
            "order", "OrderShipped", "order.shipped", order.ID,
            events.OrderShipped{ /* ... */ TrackingNo: trackingNo, Carrier: carrier})
    })
}
\end{lstlisting}
\footnotesize
\begin{itemize}\itemsep0pt
  \item \textbf{状态变更 + outbox 写} 一个事务原子 — 不会出现"状态变了但事件没发"。
  \item 下游消费者：物流系统（落 tracking 表）/ 推送服务（发"已发货"通知）/ 平台 cron（启 7d 自动确认倒计时）。
  \item 任意下游挂掉\textbf{不影响主链路}：主链路只写 outbox，publisher 异步派送。
\end{itemize}
\srcnote{internal/order/shipping.go:36--67；internal/shared/outbox/publisher.go 异步派送}
\end{frame}

% ============================================================
\section{业务边界：deck 09 不做什么（履约 gap 路线图）}
% ============================================================

\begin{frame}{业务边界：deck 09 范围之外的 6 件事}
\begin{itemize}
  \item \textbf{不真接物流商}：当前 \texttt{order.shipped} 事件里带 \texttt{tracking\_no / carrier}，但\textbf{没有 logistics\_order 子表} — 物流签收 / 拒收回写主表是路线图。
  \item \textbf{不接售后工单}：\texttt{RequestRefund / ApproveRefund / RejectRefund} 只动 order 主表的状态机，\textbf{没有 aftersales 工单表} — 退货寄回 / 验收 / 入库 / 估损扣减全部路线图。
  \item \textbf{不接评价审核}：Completed 后无评价表，"自动好评"标记 (\texttt{Auto: true}) 透传给路线图的评价服务。
  \item \textbf{不发货通知给商家}：商家工作台路线图，现版本商家凭后台 admin 接口看订单。
  \item \textbf{不接发票 / 海外清关 / 多币种}：跨境单的清关 / 汇率快照 / 关税字段路线图。
  \item \textbf{不接钱包真打款}：\texttt{ApproveRefund} 只写 \texttt{order.refunded} 事件，\texttt{TxID} 占位，下游 wallet 服务消费才真退钱；事件载荷另携带 \texttt{promo\_rule\_id / promo\_discount\_cents}，满减预算退还由 promo 侧消费同一事件异步完成（\texttt{promo\_release} 幂等台账兜重复投递）。
\end{itemize}
\keypoint{业务边界明确是\textbf{专业表达}的一部分 — 比"以为自己做了"更值钱。}
\srcnote{README.md "业务边界" 节；docs/architecture/feature-matrix.md 完整路线图}
\end{frame}

\begin{frame}{路线图 1：logistics\_order 子表 — 为什么不入主表}
\begin{itemize}
  \item 一个订单可能对应\textbf{多个物流包裹}：跨仓拆单 / 大件分箱 / 预售分批，必须 1:N。
  \item 物流子状态机（揽收 / 转运 / 派送 / 签收 / 拒收）节奏远高于主表 7 态更新频率。
  \item 物流商有 5+ 家（顺丰 / 京东 / 圆通 / EMS / 极兔），\textbf{每家 webhook JSON 不一样}，要 adapter 层抹平。
  \item 主表只保留 \texttt{order\_num}，物流字段下沉到 \texttt{logistics\_order(carrier, tracking\_no, status, ...)}。
  \item \textbf{业务粒度切分}：高频物流节点（每条物流更新）和低频订单状态（7 态切换）分两张表写，避免互相阻塞。
\end{itemize}
\srcnote{当前 internal/order/model.go:10--35 model 无 tracking 字段；路线图 repository/logistics/adapter\_xxx.go}
\end{frame}

\begin{frame}{路线图 2：aftersales 工单 — 售后完整生命周期}
\begin{tikzpicture}[node distance=6mm and 9mm,every node/.style={font=\scriptsize}]
  \node[node-box,fill=cMain!12]            (s1) {Applied\\\tiny 用户申请};
  \node[node-box,right=of s1,fill=cMain!12] (s2) {Reviewed\\\tiny 客服审核};
  \node[node-box,right=of s2,fill=cMain!12] (s3) {Returning\\\tiny 退货寄回};
  \node[node-box,right=of s3,fill=cMain!12] (s4) {Inspected\\\tiny 验收 / 估损};
  \node[node-box,right=of s4,fill=cMain!12] (s5) {Restocked\\\tiny 入库};
  \node[node-box,right=of s5,fill=cGray!18] (s6) {Refunded\\\tiny 真退款};

  \node[node-box,below=of s2,fill=cGray!18] (rj) {Rejected};
  \node[node-box,below=of s4,fill=cBad!10]  (lost){Lost\\\tiny 包裹丢失};
  \draw[arrow] (s1) -- (s2); \draw[arrow] (s2) -- (s3); \draw[arrow] (s3) -- (s4);
  \draw[arrow] (s4) -- (s5); \draw[arrow] (s5) -- (s6);
  \draw[arrow-bad,dashed] (s2) -- (rj); \draw[arrow-bad,dashed] (s3) to[bend right=18] (lost);
\end{tikzpicture}
\vspace{2mm}
\begin{itemize}
  \item 工单 6 正常态 + 2 异常态（Rejected / Lost），\textbf{独立于}主表 7 态运行。
  \item 工单只在\textbf{Restocked $\to$ Refunded} 这一步反向回写主表（推 Refunded）。
  \item \textbf{Inspected 估损}是工单和单纯退款最大的差别：钱不一定全退（货损 / 缺件按比例扣）。
\end{itemize}
\srcnote{路线图 service/aftersales/*.go；工单完结时回调 refund.ApproveRefund 完成主表收尾}
\end{frame}

\begin{frame}{路线图 3：评价表 — 挂在哪条状态边上}
\begin{tikzpicture}[node distance=6mm and 12mm,every node/.style={font=\scriptsize}]
  \node[node-box,fill=cMain!12] (wr) {WaitReceive};
  \node[node-box,right=of wr,fill=cMain!12] (cp) {Completed};
  \node[node-box,right=of cp,fill=cOK!10] (rv) {评价开窗};
  \node[node-box,below=of rv,fill=cWarn!15] (au) {Auto=true\\\tiny 系统默评};
  \node[node-box,right=of rv,fill=cMain!12] (rp) {商家回复};
  \draw[arrow] (wr) -- (cp);
  \draw[arrow] (cp) -- (rv);
  \draw[arrow,dashed] (cp) -- node[right,font=\tiny]{7d cron} (au);
  \draw[arrow] (rv) -- (rp);
\end{tikzpicture}
\vspace{2mm}
\begin{itemize}
  \item 评价窗口只在订单到 \texttt{Completed} 后开启 — 避免"没收到货就先打差评"被薅。
  \item \texttt{Auto: true}（cron 兜底）$\to$ 平台默认好评，评价服务自行决定是否计入商家评分。
  \item \texttt{Refunding} 期间评价系统\textbf{冻结}：不允许写新评价，已写的按业务规则隐藏。
\end{itemize}
\srcnote{internal/order/task.go:78--86 OrderCompletedEvent.Auto；路线图：service/review/*.go 消费此事件}
\end{frame}

\begin{frame}{现状 vs 路线图：一张图看清边界}
\begin{tikzpicture}[node distance=4mm and 10mm,every node/.style={font=\scriptsize}]
  \node[node-box,node-ok] (a) {订单 7 态机\\\tiny 已实现};
  \node[node-box,node-ok,right=of a] (b) {Outbox 事件\\\tiny 已实现};
  \node[node-box,node-ok,right=of b] (c) {RMQ+Cron 关单\\\tiny 已实现（已加固）};
  \node[node-box,node-ok,below=of a] (d) {7d 自动确认\\\tiny 已实现};
  \node[node-box,node-ok,right=of d] (e) {Saga 释放预占\\\tiny 已实现};
  \node[node-box,node-ok,right=of e] (f) {旧别名清退\\\tiny 已完成（已删）};
  \node[node-box,node-warn,below=of d] (g) {logistics 子表\\\tiny 路线图};
  \node[node-box,node-warn,right=of g] (h) {aftersales 工单\\\tiny 路线图};
  \node[node-box,node-warn,right=of h] (i) {评价 / 默评\\\tiny 路线图};
  \node[node-box,node-warn,below=of g] (j) {商家工作台\\\tiny 路线图};
  \node[node-box,node-warn,right=of j] (k) {钱包真打款\\\tiny 路线图};
  \node[node-box,node-warn,right=of k] (l) {多币种 / 跨境\\\tiny 路线图};
\end{tikzpicture}
\vspace{2mm}
\begin{itemize}
  \item 上 6 项已落地（其中关单 cron 经历过一次表达式陷阱并已加固），下 6 项是路线图。
  \item gomall MVP 聚焦\textbf{交易闭环} — 状态机推到 7 态，下游服务由 outbox 解耦后续逐步落地。
\end{itemize}
\srcnote{README.md "技术覆盖"+"业务边界" 双表；docs/architecture/feature-matrix.md}
\end{frame}

% ============================================================
\section{兼容性：3 $\to$ 7 零停机迁移}
% ============================================================

\begin{frame}{零停机迁移：DB 不动 + 别名兼容}
\begin{itemize}
  \item \textbf{硬约束}：DB 有数百万存量 order，\texttt{type} 字段是 tinyint，改值需 ALTER + UPDATE 全表扫，\textbf{停机不可接受}。
  \item \textbf{老旧客户端}（小程序 / Android / 历史 API）还在用旧字段名 \texttt{UnPaid / Paid / Cancelled} 拼前端字典。
  \item \textbf{迁移策略}：
  \begin{itemize}
    \item DB 数值\textbf{完全不动}：1 = WaitPay（旧名 UnPaid），语义对齐。
    \item 迁移期旧名以 \texttt{Deprecated} 常量别名共存，老调用方编译不破。
    \item 旧前端字典 \texttt{OrderTypeMap} 迁移期并存，转译同一组 int。
  \end{itemize}
  \item \textbf{编译期校验}：\texttt{// Deprecated:} 让 staticcheck / golangci-lint warn 出来，新代码不写老名。
  \item \textbf{收尾}：调用方全部切换后，别名块与 \texttt{OrderTypeMap} \textbf{已整体删除}，现行字典只剩 \texttt{OrderStateMap}。
\end{itemize}
\keypoint{常量层别名 + 编译期校验 + 按期清退 = \textbf{零停机 / 零回归}的迁移。}
\srcnote{consts/order.go:6--28 现行常量与 OrderStateMap；Deprecated 别名与 OrderTypeMap 已随迁移收尾删除}
\end{frame}

% ============================================================
\appendix
% ============================================================

\begin{frame}{业务 Q\&A（上）}
\textbf{Q1}：为什么 30min 关单不能由商家自定义？\\
\hspace{4mm}\small A：平台规则统一 = 用户预期统一 = 客服话术统一。商家自定义会导致客诉口径混乱（路线图：仅大件 / 预售类可商家配）。\\
\hspace{4mm}\srcnote{repository/rabbitmq/order\_delay.go:22 默认 30min 常量}

\vspace{1mm}
\textbf{Q2}：当初的 cron 表达式陷阱为什么没造成线上事故？现在怎么收口的？\\
\hspace{4mm}\small A：靠 \texttt{CloseOrderWithCheck} 的 \texttt{WHERE type=1} 条件 UPDATE 兜底幂等 + 释放 reserved 是 Lua 幂等 + 关单不可逆，所以误写期间最坏只是 60× 无效 DB 空跑、不丢单不超卖；现行已改用 \texttt{@every 5m} 从写法上消除秒位歧义。\\
\hspace{4mm}\srcnote{initialize/cron.go:21 现行 @every 5m + internal/order/cancel.go:32 + internal/order/repo.go:198}

\vspace{1mm}
\textbf{Q3}：用户付了 3 天商家没发货，怎么自动处理？\\
\hspace{4mm}\small A：当前靠客服 trigger \texttt{RequestRefund}（from=WaitShip 合法）；路线图加 cron 扫 WaitShip+48h 自动 Refunding。\\
\hspace{4mm}\srcnote{internal/refund/service.go:37--41 from 含 WaitShip}

\vspace{1mm}
\textbf{Q4}：7d 自动确认收货对商家公平吗？\\
\hspace{4mm}\small A：业界共识（淘宝实物 15d / 京东 7d / 拼多多 7d）。Completed 后\textbf{仍可发起退款} $\to$ Refunding，商家利益不会被一次性结算锁死。\\
\hspace{4mm}\srcnote{internal/order/task.go:17 autoConfirmReceiveDays=7；internal/order/state.go:25 Completed 出边}

\vspace{1mm}
\textbf{Q5}：用户能撤回"已确认收货"吗？\\
\hspace{4mm}\small A：不能（状态机不允许 Completed $\to$ WaitReceive 回退）。但走 \texttt{Completed $\to$ Refunding} 售后路径仍可退款。\\
\hspace{4mm}\srcnote{internal/order/state.go:21--28 转换表}
\end{frame}

\begin{frame}{业务 Q\&A（下）}
\textbf{Q6}：商家发货 24h 后又想撤销发货，能吗？\\
\hspace{4mm}\small A：不能（WaitReceive $\to$ WaitShip 不存在）。商家联系客服走\textbf{用户全额退款}流程。\\
\hspace{4mm}\srcnote{internal/order/state.go:24 WaitReceive 出边只到 Completed / Refunding}

\vspace{1mm}
\textbf{Q7}：物流单号不入表，下游怎么知道？\\
\hspace{4mm}\small A：\texttt{order.shipped} 事件 payload 带 \texttt{tracking\_no / carrier}，下游物流服务自行落表；标准 outbox 解耦。\\
\hspace{4mm}\srcnote{internal/order/shipping.go:36--67 + events/events.go OrderShipped}

\vspace{1mm}
\textbf{Q8}：退款多久能到账？\\
\hspace{4mm}\small A：状态机本身瞬时切换（\texttt{Refunding $\to$ Refunded}）；真打款 1-15d 由下游钱包 / 支付通道决定（路线图）。\\
\hspace{4mm}\srcnote{events/events.go OrderRefundedEvent.TxID 占位；wallet 服务消费事件}

\vspace{1mm}
\textbf{Q9}：商家在用户申请退款期间能继续发货吗？\\
\hspace{4mm}\small A：不能。\texttt{ShipOrder} DAO 的 \texttt{WHERE type=WaitShip} 不会匹配 Refunding 行，RowsAffected=0 直接拒。\\
\hspace{4mm}\srcnote{internal/order/repo.go:215--223}

\vspace{1mm}
\textbf{Q10}：3 态 $\to$ 7 态要发停机公告吗？\\
\hspace{4mm}\small A：不用。DB 数值 0 变更，迁移期常量层 alias 兜住编译，调用方切完后别名已删。\\
\hspace{4mm}\srcnote{consts/order.go:6--16 现行常量}

\vspace{1mm}
\textbf{Q11}：deck 11 讲的 outbox 和本 deck 矛盾吗？\\
\hspace{4mm}\small A：不矛盾。deck 11 讲 outbox 模式本身，本 deck 是\textbf{应用方}：每个状态切换都走 outbox 解耦下游。\\
\hspace{4mm}\srcnote{internal/shared/outbox/publisher.go；本 deck 5 处 \texttt{NewOutboxDaoByDB(tx).Insert}}

\vspace{1mm}
\textbf{Q12}：订单链路压测数字怎么对应业务承诺？\\
\hspace{4mm}\small A：\texttt{/orders/create} 50K RPS / p95 2.33ms 对应 P0 SLO（99.95\% / p99<500ms）留 \textbf{200× 余量}；\texttt{/orders/list} 58K RPS 对应 P1 SLO。\\
\hspace{4mm}\srcnote{stressTest/REPORT.md 第 4 节幂等下单；第 2 节订单列表}
\end{frame}

\begin{frame}{代码位置一览}
\begin{description}
  \item[状态机表 + 工具函数]   \srcnote{internal/order/state.go:21--61}
  \item[状态常量（旧别名已删）] \srcnote{consts/order.go:6--16}
  \item[OrderStateMap 中文名]  \srcnote{consts/order.go:19--28}
  \item[ShipOrder service]    \srcnote{internal/order/shipping.go:36--67}
  \item[ConfirmReceive service] \srcnote{internal/order/shipping.go:73--111}
  \item[RequestRefund service]  \srcnote{internal/refund/service.go:47--87}
  \item[ApproveRefund service]  \srcnote{internal/refund/service.go:94--135}
  \item[RejectRefund service]   \srcnote{internal/refund/service.go:140--170}
  \item[CancelUnpaidOrder 共享入口] \srcnote{internal/order/cancel.go:19--67}
  \item[满减预算退还消费侧（异步）]  \srcnote{internal/promo/consumer.go:25--48；initialize/promo.go}
  \item[CloseOrderWithCheck DAO]    \srcnote{internal/order/repo.go:198--209}
  \item[ShipOrder DAO]              \srcnote{internal/order/repo.go:215--223}
  \item[ConfirmReceive DAO]         \srcnote{internal/order/repo.go:227--235}
  \item[RequestRefund DAO]          \srcnote{internal/order/repo.go:239--250}
  \item[GetTimeoutWaitReceive DAO]  \srcnote{internal/order/repo.go:277--290}
  \item[Cron 自动确认收货]          \srcnote{internal/order/task.go:61--94}
  \item[Cron 注册 (现行 @every 5m)]  \srcnote{initialize/cron.go:19--31}
  \item[RMQ 延迟队列 30min]         \srcnote{repository/rabbitmq/order\_delay.go:22 + 59--72}
  \item[业务码 70001 / 70002 / 60002] \srcnote{pkg/e/code.go:54 + 57--58}
\end{description}
\end{frame}

\end{document}
